\sect{Эксперименты}{experiments}
%
Эксперименты направлены на проверку трёх гипотез: (i) гибридная схема
HMM+OLTW сохраняет точность выравнивания при наличии rubato и редких
ошибок исполнителя; (ii) свёрточный фронтенд снижает ошибку
выравнивания на полифоническом материале по сравнению с гауссовой
эмиссией; (iii) RL-модуль опережающего предсказания темпа уменьшает
эффективную задержку аккомпанемента без ухудшения точности.

Базовые конфигурации, с которыми проводилось сравнение, обозначены
далее как \textit{HMM} (только скрытая марковская модель с гауссовой
эмиссией и фиксированной полосной матрицей переходов), \textit{HMM-CNN}
(тот же тракт, но с обучаемой эмиссией из~\Secref{cnn-corner}) и
\textit{HMM+OLTW+RL} (полная предложенная схема). Все три конфигурации
работают в общем конвейере (\Secref{realtime}) с одинаковыми
параметрами очистки и одинаковым рендерером, что позволяет
сосредоточить сравнение на алгоритмических различиях.

Метрики качества подобраны в соответствии с принятой в литературе
практикой~\cite{Nakamura2015,Henkel2025}. Средняя абсолютная ошибка
выравнивания MAE измеряется в номерах нот и характеризует, насколько
далеко предсказанная позиция отличается от истинной по партитуре.
Метрика F1-score рассчитывается как гармоническое среднее точности и
полноты для onset-событий, выровненных с допуском $50$\,мс. Метрика
ReSync~--- доля исполнения, на которой потребовалась повторная
синхронизация. Эффективная задержка $L_{\text{eff}}$ измеряется
как медиана временного сдвига между отрисованным оркестровым сэмплом
и соответствующим onset солиста.

Сначала рассмотрим зависимость задержки аккомпанемента от горизонта
предсказания темпа. На \Figref{tempo-latency} видно, что без
опережающего модуля эффективная задержка монотонно возрастает с
ростом горизонта (поскольку рендерер вынужден буферизовать всё
больший интервал), тогда как для гибридной схемы HMM+OLTW+RL зависимость
немонотонна и достигает минимума в области $250$--$350$\,мс. Это
согласуется с проектным значением горизонта $\gamma=0{,}95$.

\par\addvspace{2pt}
{\centering
\refstepcounter{figure}\figlab{tempo-latency}%
\resizebox{\columnwidth}{!}{% =====================================================================
% TikZ figure: tempo-anticipation horizon vs effective accompaniment
% latency (BLACK & WHITE).  Used in sections/08-experiments.tex.
% Axis labels are placed near the centre of each axis (not at the
% arrow tips) so they do not collide with the arrowheads.
% =====================================================================
\begin{tikzpicture}[x=0.9cm,y=0.075cm]
  \draw[->,thick] (0,0) -- (8.2,0);
  \draw[->,thick] (0,0) -- (0,62);

  \foreach \x/\labl in {1/100,2/150,3/200,4/250,5/300,6/400,7/500}
    \draw (\x,0) -- (\x,-1.5) node[below,font=\scriptsize] {\labl};
  \foreach \y in {10,20,30,40,50,60}
    \draw (0,\y) -- (-0.12,\y) node[left,font=\scriptsize] {\y};

  \foreach \y in {10,20,30,40,50}
    \draw[dotted,black!50] (0,\y) -- (7.2,\y);

  \draw[thick]
    plot[smooth] coordinates
      {(1,44) (2,41) (3,37) (4,34) (5,31) (6,28) (7,26)};
  \draw[thick,dashed]
    plot[smooth] coordinates
      {(1,23) (2,21) (3,20) (4,19) (5,19) (6,20) (7,21)};

  \node[font=\scriptsize,anchor=west] at (4.4,46) {базовый трекер};
  \node[font=\scriptsize,anchor=west] at (4.4,16) {HMM+OLTW+RL};

  \node[font=\scriptsize,below=4mm] at (4,-3) {Горизонт, мс};
  \node[font=\scriptsize,rotate=90] at (-1.0,32) {Задержка, мс};
\end{tikzpicture}
}\par
\vspace{2pt}
\footnotesize
\figurename~\thefigure. Зависимость эффективной задержки
аккомпанемента от горизонта предсказания темпа для базового
трекера и предложенной гибридной схемы.\par}
\addvspace{4pt}

Далее измерена устойчивость трекера при увеличении вариативности
rubato (\Figref{error-vs-rubato}). Базовый HMM начинает терять точность
выравнивания при $\sigma_\tau \ge 6$\,\%, тогда как HMM+OLTW+RL
сохраняет приемлемую ошибку вплоть до $\sigma_\tau\approx 12$\,\%.
Это объясняется двумя факторами: возможностью переключения на OLTW в
зонах низкой уверенности HMM (\Secref{realtime}) и опережающим
управлением темпом (\Secref{rl}), которое сглаживает резкие
ускорения.

\par\addvspace{2pt}
{\centering
\refstepcounter{figure}\figlab{error-vs-rubato}%
\resizebox{\columnwidth}{!}{% =====================================================================
% TikZ figure: rubato strength vs alignment error (BLACK & WHITE).
% Used in sections/08-experiments.tex.
% Axis labels are placed near the centre of each axis (not at the
% arrow tips) so they do not collide with the arrowheads.
% =====================================================================
\begin{tikzpicture}[x=0.95cm,y=0.10cm]
  \draw[->,thick] (0,0) -- (8.0,0);
  \draw[->,thick] (0,0) -- (0,42);

  \foreach \x/\labl in {1/2,2/4,3/6,4/8,5/10,6/12,7/14}
    \draw (\x,0) -- (\x,-1.1) node[below,font=\scriptsize] {\labl};
  \foreach \y in {5,10,15,20,25,30,35,40}
    \draw (0,\y) -- (-0.11,\y) node[left,font=\scriptsize] {\y};

  \foreach \y in {10,20,30}
    \draw[dotted,black!50] (0,\y) -- (7.2,\y);

  \draw[thick]
    plot[smooth] coordinates
      {(1,7) (2,10) (3,14) (4,19) (5,25) (6,31) (7,36)};
  \draw[thick,dashed]
    plot[smooth] coordinates
      {(1,6) (2,8) (3,11) (4,15) (5,19) (6,24) (7,28)};

  \node[font=\scriptsize,anchor=west] at (4.4,33) {HMM};
  \node[font=\scriptsize,anchor=west] at (4.4,22) {HMM+OLTW+RL};

  \node[font=\scriptsize,below=4mm] at (4,-3) {Сила rubato, \%};
  \node[font=\scriptsize,rotate=90] at (-1.1,21) {Ошибка, мс};
\end{tikzpicture}
}\par
\vspace{2pt}
\footnotesize
\figurename~\thefigure. Зависимость ошибки выравнивания от силы
rubato. Базовая схема HMM теряет точность раньше, чем гибридная
схема HMM+OLTW+RL.\par}
\addvspace{4pt}

Сводные количественные результаты по всем трём конфигурациям
приведены в \Tabref{main-results}. Гибридная схема HMM+OLTW+RL
превосходит базовый HMM на тестовой части CU-Concerto-2026 по всем
четырём метрикам: MAE снижается с $0{,}19$ до $0{,}11$ ноты, F1
возрастает с $0{,}947$ до $0{,}979$, ReSync падает в $2{,}5$ раза, а
эффективная задержка $L_{\text{eff}}$ опускается с $34$ до $19$\,мс.
Замена гауссовой эмиссии на CNN (HMM-CNN) даёт промежуточный
результат, что подтверждает значимость каждого из двух нововведений.

\par\addvspace{2pt}
{\centering
\refstepcounter{table}\tablab{main-results}%
\footnotesize
\begin{tabular}{@{}lcccc@{}}
  \toprule
  Конфигурация       & MAE, нот & F1     & ReSync, \% & $L_{\text{eff}}$, мс \\
  \midrule
  HMM                & $0{,}19$ & $0{,}947$ & $5{,}3$  & $34$ \\
  HMM-CNN            & $0{,}14$ & $0{,}968$ & $3{,}1$  & $30$ \\
  HMM+OLTW+RL        & $0{,}11$ & $0{,}979$ & $2{,}1$  & $19$ \\
  \bottomrule
\end{tabular}\par
\vspace{2pt}
\tablename~\thetable. Сводные количественные результаты на
тестовой части датасета CU-Concerto-2026. Меньшие значения MAE,
ReSync и $L_{\text{eff}}$~--- лучше; большее значение F1~--- лучше.\par}
\addvspace{4pt}

Результаты по отдельным сценариям исполнения, включающим
естественное rubato, пропуски нот, орнаменты и педальные удержания,
приведены в \Tabref{scenarios}. Падение F1 у предложенной схемы
наблюдается только в самых тяжёлых режимах (комбинированный стресс-тест
и низкое отношение сигнал/шум аудио), но даже в них значение
остаётся выше $0{,}93$.

\par\addvspace{2pt}
{\centering
\refstepcounter{table}\tablab{scenarios}%
\footnotesize
\begin{tabular}{@{}p{0.45\linewidth}cccc@{}}
  \toprule
  Сценарий исполнения                          & $\sigma_\tau$, \% & MAE & F1 & ReSync, \% \\
  \midrule
  Ровный темп, без ошибок                      & 2  & $0{,}07$ & $0{,}992$ & $0{,}0$ \\
  Умеренное rubato, без пропусков              & 6  & $0{,}09$ & $0{,}988$ & $0{,}0$ \\
  Сильное rubato в каденции                    & 11 & $0{,}12$ & $0{,}981$ & $0{,}7$ \\
  Случайные ускорения на переходах             & 8  & $0{,}11$ & $0{,}984$ & $0{,}5$ \\
  Резкое замедление перед кульминацией         & 10 & $0{,}13$ & $0{,}979$ & $0{,}9$ \\
  Единичный пропуск ноты                       & 4  & $0{,}10$ & $0{,}986$ & $1{,}1$ \\
  Серия из двух пропусков                      & 5  & $0{,}14$ & $0{,}975$ & $2{,}4$ \\
  Ложное повторение короткого мотива           & 7  & $0{,}16$ & $0{,}969$ & $3{,}1$ \\
  Орнамент с лишними быстрыми нотами           & 6  & $0{,}15$ & $0{,}972$ & $2{,}7$ \\
  Переход с педалью и размытым onset           & 7  & $0{,}18$ & $0{,}964$ & $4{,}8$ \\
  Низкое отношение сигнал/шум (аудио)          & 8  & $0{,}21$ & $0{,}951$ & $6{,}5$ \\
  Комбинированный стресс-тест                  & 12 & $0{,}27$ & $0{,}932$ & $9{,}8$ \\
  \bottomrule
\end{tabular}\par
\vspace{2pt}
\tablename~\thetable. Метрики по отдельным сценариям исполнения
для предложенной схемы HMM+OLTW+RL на CU-Concerto-2026.\par}
\addvspace{4pt}

Дополнительно проведена субъективная оценка с участием шести
студентов-консерваторов, не участвовавших в записи обучающего
датасета. Каждому испытуемому предлагалось сыграть одну из частей
концерта Чайковского с виртуальным аккомпанементом по двум
конфигурациям (HMM и HMM+OLTW+RL) в случайном порядке без указания,
какая из них активна. По завершении испытуемый оценивал три аспекта
по пятибалльной шкале: естественность темпового следования,
ощущение «живой реакции» оркестра и общее удобство репетиции.
Средние оценки HMM+OLTW+RL по всем трём аспектам выше базовой
конфигурации на $1{,}1$ балла, что согласуется с количественными
результатами и подтверждает практическую значимость
предложенной схемы.
